\documentclass[aspectratio=169]{beamer}
\input{header}
\coursecode{JKSSB}

\title{Output Devices}
\subtitle{Lesson 14 --- Displays, Audio and Hard vs Soft Copy}
\author{JKSSB Computer Awareness}
\date{\today}

\begin{document}
\frame{\titlepage}

\begin{frame}{What Is an Output Device?}
  \begin{itemize}
    \item An output device \key{delivers the result} of processing to the user
      or to another device.
    \item It converts machine-readable data into a form a human can
      \key{see, hear, or read}.
    \item Common examples: monitor, projector, speaker, headphones, printer,
      and plotter.
    \item \muted{Output is the last step of the input--process--output cycle.}
  \end{itemize}
  \begin{block}{Definition}
    An output device receives processed data from the CPU and presents it to
    the user in a readable form.
  \end{block}
\end{frame}

\begin{frame}{Classifying Output Devices}
  \begin{center}
    \begin{tabular}{p{0.26\textwidth}p{0.60\textwidth}}
      \toprule
      \textbf{Category} & \textbf{Typical devices} \\
      \midrule
      Visual / soft copy & Monitor, projector \\
      Audio & Speaker, headphones \\
      Printed / hard copy & Printer, plotter \\
      \bottomrule
    \end{tabular}
  \end{center}
  \vspace{0.2cm}
  \muted{A monitor and a projector show output on a screen; a printer and a
  plotter put it on paper.}
\end{frame}

\begin{frame}{Soft Copy vs Hard Copy}
  \begin{center}
    \begin{tabular}{p{0.18\textwidth}p{0.34\textwidth}p{0.34\textwidth}}
      \toprule
      \textbf{Point} & \textbf{Soft copy} & \textbf{Hard copy} \\
      \midrule
      Form & Electronic, on a screen & Physical, on paper \\
      Examples & Monitor, projector & Printer, plotter \\
      Editable & Easily edited & Fixed once printed \\
      \bottomrule
    \end{tabular}
  \end{center}
  \vspace{0.25cm}
  \begin{alertblock}{Trap alert}
    \key{Soft copy} is displayed on a screen and is not printed;
    \key{hard copy} is printed on paper. Do not swap the two.
  \end{alertblock}
\end{frame}

\begin{frame}{Monitors}
  \begin{itemize}
    \item A \key{monitor} is the main visual output device of a computer.
    \item It displays text, images, and video as \key{soft copy}.
    \item Screen size is measured \key{diagonally in inches}; picture
      sharpness is measured in \key{pixels} (resolution).
    \item Common technologies are \key{CRT}, \key{LCD}, \key{LED}, and
      \key{OLED}.
  \end{itemize}
  \begin{block}{Key idea}
    A monitor shows output to the user; it does not take input (a touchscreen
    is the exception, because it also senses touch).
  \end{block}
\end{frame}

\begin{frame}{CRT vs Flat-Panel Displays}
  \begin{center}
    \begin{tabular}{p{0.20\textwidth}p{0.34\textwidth}p{0.34\textwidth}}
      \toprule
      \textbf{Point} & \textbf{CRT} & \textbf{Flat panel} \\
      \midrule
      Full form & Cathode Ray Tube & LCD / LED / OLED \\
      Shape & Bulky, deep & Thin, flat, light \\
      Technology & Electron beam in a vacuum tube & Liquid crystal or self-emitting pixels \\
      \bottomrule
    \end{tabular}
  \end{center}
  \vspace{0.25cm}
  \begin{alertblock}{Trap alert}
    A \key{CRT} (Cathode Ray Tube) is the older, bulky display. Modern
    monitors are \key{flat panels}.
  \end{alertblock}
\end{frame}

\begin{frame}{LCD, LED, and OLED}
  \begin{center}
    \begin{tabular}{p{0.14\textwidth}p{0.42\textwidth}p{0.24\textwidth}}
      \toprule
      \textbf{Type} & \textbf{What it means} & \textbf{Backlight} \\
      \midrule
      LCD & Liquid Crystal Display & Needs a backlight \\
      LED & Light Emitting Diode display (an LCD panel lit by LEDs) & LED backlight \\
      OLED & Organic Light Emitting Diode & Self-emitting, no backlight \\
      \bottomrule
    \end{tabular}
  \end{center}
  \vspace{0.2cm}
  \muted{In an LED display the LEDs provide the light; in an OLED display
  each pixel makes its own light.}
\end{frame}

\begin{frame}{Projector}
  \begin{itemize}
    \item A \key{projector} throws a large image onto a screen or wall.
    \item It is used in classrooms, meetings, and cinemas.
    \item The projected image is \key{soft copy}, because it is shown, not
      printed.
    \item \muted{A projector is an output device; it does not capture
      anything.}
  \end{itemize}
\end{frame}

\begin{frame}{Audio Output: Speakers and Headphones}
  \begin{itemize}
    \item A \key{speaker} converts an electrical signal into sound for a
      group of listeners.
    \item \key{Headphones} produce sound for a \key{single listener} and are
      worn on or over the ears.
    \item Both are \key{audio output} devices.
  \end{itemize}
  \begin{block}{Pair to remember}
    A \key{microphone} is audio \key{input}; a speaker or headphone is audio
    \key{output}.
  \end{block}
\end{frame}

\begin{frame}{Printers: Overview}
  \begin{itemize}
    \item A \key{printer} produces \key{hard copy} on paper.
    \item Printers fall into two broad classes: \key{impact} and
      \key{non-impact}.
    \item Common types: \key{dot-matrix} (impact); \key{inkjet} and
      \key{laser} (non-impact).
    \item Lasers are common in offices; inkjets are common in homes.
    \item \muted{A printer is an output device (TRAP-12).}
  \end{itemize}
\end{frame}

\begin{frame}{Impact vs Non-Impact Printers}
  \begin{center}
    \begin{tabular}{p{0.18\textwidth}p{0.34\textwidth}p{0.34\textwidth}}
      \toprule
      \textbf{Point} & \textbf{Impact} & \textbf{Non-impact} \\
      \midrule
      Method & Strikes a ribbon and the paper & No striking; uses ink, toner, or heat \\
      Examples & Dot-matrix & Inkjet, laser \\
      Noise & Noisier & Quieter \\
      Graphics quality & Lower & Higher \\
      \bottomrule
    \end{tabular}
  \end{center}
  \vspace{0.2cm}
  \begin{alertblock}{Trap alert}
    A \key{dot-matrix} printer is an \key{impact} printer; \key{inkjet} and
    \key{laser} printers are \key{non-impact}.
  \end{alertblock}
\end{frame}

\begin{frame}{Plotter (PYQ-29)}
  \begin{itemize}
    \item A \key{plotter} draws large, precise \key{line drawings} such as
      engineering and architectural designs.
    \item It draws using a pen or a vector-based method, on paper larger than
      a normal printer's.
    \item \muted{A plotter is an output device, not an input device
      (TRAP-12).}
  \end{itemize}
  \begin{exampleblock}{PYQ-29}
    Which output device is used to produce \key{large engineering drawings}?
    Answer: a \key{plotter}.
  \end{exampleblock}
\end{frame}

\begin{frame}{Multifunction Printer (MFP)}
  \begin{itemize}
    \item \key{MFP} stands for \key{Multifunction Printer}; it is also called
      an \key{all-in-one} device.
    \item It combines \key{printing, scanning, copying}, and often
      \key{faxing} in a single unit.
    \item Its scanning part is \key{input}; its printing part is
      \key{output}.
    \item Common in offices and small businesses.
  \end{itemize}
  \begin{block}{Keep it straight}
    An MFP both takes input (when scanning) and gives output (when printing).
  \end{block}
\end{frame}

\begin{frame}{Trap Alert: Output or Input? (TRAP-12)}
  \begin{itemize}
    \item \key{Output}: monitor, projector, speaker, headphones, printer,
      plotter, drum printer, sound card.
    \item \key{Input}: keyboard, mouse, scanner, light pen, digitizer tablet,
      punch-card reader, microphone.
    \item \muted{A plotter and a printer give output; a scanner, light pen,
      and digitizer give input.}
  \end{itemize}
  \begin{alertblock}{Trap alert}
    A plotter, a printer, and a sound card are \key{output}. A scanner and a
    digitizer are \key{input}, even when grouped together in one option list.
  \end{alertblock}
\end{frame}

\begin{frame}{Near-Neighbours to Keep Apart}
  \begin{center}
    \begin{tabular}{p{0.30\textwidth}p{0.60\textwidth}}
      \toprule
      \textbf{Pair} & \textbf{The distinction to state} \\
      \midrule
      Printer vs plotter & Printer: standard pages of text/graphics. Plotter: large, precise line drawings. \\
      Monitor vs projector & Monitor: image on its own screen. Projector: image thrown onto a screen or wall. \\
      Speaker vs microphone & Speaker: audio output. Microphone: audio input. \\
      Soft copy vs hard copy & Soft copy: shown on a screen. Hard copy: printed on paper. \\
      \bottomrule
    \end{tabular}
  \end{center}
\end{frame}

\begin{frame}{Lesson 14: Recall}
  \begin{itemize}
    \item An output device presents processed data in a form the user can see,
      hear, or read.
    \item Soft copy is on a screen (monitor, projector); hard copy is on
      paper (printer, plotter).
    \item Monitor types: CRT (bulky, Cathode Ray Tube); flat panels LCD
      (backlight), LED (LED backlight), OLED (self-emitting).
    \item Audio output: speakers and headphones. A microphone is audio input.
    \item Printers: impact (dot-matrix) vs non-impact (inkjet, laser). A
      plotter draws large engineering drawings (PYQ-29).
    \item MFP = Multifunction Printer (print, scan, copy, often fax).
    \item TRAP-12: plotter, printer, and sound card are output; scanner,
      light pen, and digitizer are input.
  \end{itemize}
\end{frame}
\end{document}
